\chapter{Interfaccia web: visualizzazione e raccolta dati}
\label{cap:webui}

\begin{quote}\small
Questo capitolo risponde all'obiettivo O5: l'ESP32 pubblica i dati, un sito web li
visualizza in tempo reale, la stessa applicazione li salva con statistiche ed
esporta in CSV per l'analisi numerica. Alla data di questa stesura il progetto
architetturale e di dettaglio è completo; l'implementazione è in corso.
\end{quote}

\section{Requisiti e scelta architetturale}
\label{sec:webui-requisiti}

\subsection{Requisiti}

Dall'obiettivo O5 si ricavano cinque requisiti funzionali e due non funzionali.

\begin{description}[leftmargin=1.4em,style=nextline,itemsep=2pt]
  \item[R1 --- Pubblicazione] l'ESP32 rende disponibili i dati del sensore in rete.
  \item[R2 --- Tempo reale] la visualizzazione segue il sensore senza ritardi
    percepibili, alla stessa cadenza dell'acquisizione (5~Hz).
  \item[R3 --- Statistiche] la stessa applicazione calcola e mostra statistiche di
    sessione.
  \item[R4 --- Esportazione] i dati sono esportabili in CSV per l'analisi in Excel.
  \item[R5 --- Vitalità] l'interfaccia mostra anche l'indice di vitalità
    (\cref{cap:vitalita}).
  \item[N1 --- Funzionamento offline] vincolo derivato dal contesto UPRISE: in
    emergenza sismica non c'è connettività.
  \item[N2 --- Coerenza dei dati] il CSV esportato dal browser deve avere lo stesso
    formato di quello prodotto dalla catena seriale, in modo che esista
    \emph{un solo} formato dati in tutta la tesi.
\end{description}

\subsection{Opzioni considerate}

\begin{table*}[htbp]
\centering
\caption{Opzioni architetturali valutate per la web UI.}
\label{tab:webui-opzioni}
\small
\begin{tabularx}{\linewidth}{@{}lXX@{}}
\toprule
\textbf{Architettura} & \textbf{Vantaggi} & \textbf{Svantaggi} \\
\midrule
\textbf{A. Sito servito dall'ESP32} (web server e WebSocket a bordo) &
  nessuna infrastruttura, funziona offline, dimostrazione autonoma, coerente con lo
  scenario di emergenza &
  RAM e flash limitate, massimo 2--3 client simultanei \\
B. ESP32 $\to$ MQTT $\to$ server con database &
  scalabile, storico illimitato, più vicino a una piattaforma reale &
  richiede broker e server sempre accesi: infrastruttura eccessiva per una tesi di
  sensing \\
C. Piattaforma cloud (ThingSpeak, Blynk) &
  pronta all'uso &
  non funziona offline, cadenza di campionamento limitata ($\sim$15~s), dati fuori
  controllo \\
\bottomrule
\end{tabularx}
\end{table*}

\subsection{Scelta: tutto sull'ESP32}

È stata scelta l'opzione~A, per tre motivi.

\begin{enumerate}[leftmargin=1.8em,itemsep=3pt]
  \item \textbf{Coerenza con il contesto}: un nodo autonomo che serve la propria
        dashboard è la miniatura concettuale della piattaforma di monitoraggio del
        progetto, con le sue due modalità di funzionamento. L'opzione~C viola
        direttamente il requisito N1.
  \item \textbf{Il volume dei dati è piccolo}: circa venti valori a 5~Hz,
        dell'ordine di 2~kB/s, ben dentro le capacità dell'ESP32.
  \item \textbf{La persistenza a lungo termine non serve a bordo}: lo storico si
        accumula \emph{nel browser}, che ha memoria abbondante, ed è il browser a
        produrre il CSV. Questo elimina il bisogno di un database senza perdere
        funzionalità.
\end{enumerate}

\paragraph{Modalità WiFi doppia.} Il firmware supporta due modalità selezionabili:
\emph{Station}, in cui l'ESP32 si collega a una rete esistente (comodo in fase di
sviluppo), e \emph{Access Point}, in cui crea la propria rete e serve la pagina a
chi si collega. Se la connessione in modalità Station non riesce entro 15~s, parte
automaticamente l'Access Point. È la modalità che rende la dimostrazione possibile
in qualunque luogo, e quella che corrisponde allo scenario di emergenza.

\paragraph{Piano di riserva.} È stato progettato anche il percorso alternativo
(opzione~B: ESP32 $\to$ MQTT $\to$ backend FastAPI con SQLite $\to$ browser) nel
caso in cui il relatore intenda una piattaforma centralizzata anziché il nodo
autonomo. Il frontend è condiviso per circa l'85\,\% fra le due architetture, quindi
il cambio di rotta costerebbe due o tre giornate di lavoro. La scelta di riferimento
resta l'opzione~A.

\section{Progetto}
\label{sec:webui-progetto}

\subsection{Stack}

Lato firmware: \texttt{ESPAsyncWebServer} per servire la pagina, WebSocket per il
flusso dati, LittleFS per ospitare i file statici. Tutto risiede sull'ESP32 e nulla
è caricato da rete esterna, in ottemperanza al requisito N1.

Lato browser: pagina singola in HTML, CSS e JavaScript, senza framework e senza
dipendenze da CDN --- coerentemente con il vincolo di funzionamento offline.

\subsection{Protocollo dati}

Il messaggio WebSocket riprende le stesse grandezze del CSV, con l'aggiunta dei campi
calcolati a bordo:

\begin{lstlisting}[caption={Messaggio WebSocket inviato a 5~Hz.},
label={lst:ws},basicstyle=\ttfamily\scriptsize]
{
  "t": 123456,                        // millis() dell'ESP32
  "presence": 1,                      // radar_presence
  "moving": 1, "still": 0,            // stato del bersaglio
  "mdist": 145, "sdist": 0,           // distanze in cm
  "menergy": 72, "senergy": 0,        // energia del bersaglio 0-100
  "pir": 1,                           // sensore PIR (confronto dal vivo)
  "gates_m": [0,5,60,10,0,0,0,0,0],   // energia per-gate, movimento
  "gates_s": [0,0,45,8,0,0,0,0,0],    // energia per-gate, stazionario
  "vitality": 54,                     // indice di vitalita 0-100
  "vitality_class": "attivo"          // nessun_segno | vitalita_bassa |
                                      // moderato | attivo
}
\end{lstlisting}

La cadenza di invio è di \textbf{5~Hz}, identica al campionamento: nessun
sottocampionamento, così che il CSV esportato dal browser sia equivalente a quello
prodotto dalla catena seriale (requisito N2).

Il canale WebSocket gestisce la riconnessione automatica con backoff: se la
connessione cade, la pagina segnala lo stato e riprende senza perdere le statistiche
già accumulate.

\subsection{Struttura della pagina}

La pagina è organizzata in quattro aree, schematizzate in \cref{fig:webui-layout}.

\begin{figure*}[htbp]
\centering
\begin{minipage}{0.95\linewidth}
\ttfamily\scriptsize
\begin{verbatim}
+-------------------------------------------------------------+
|  HEADER   * PRESENZA RILEVATA     [dist. 1.45 m]  [REC]     |
+--------------------------+----------------------------------+
|  A. STATO LIVE           |  B. INDICE DI VITALITA'          |
|  radar: MOVING @ 145 cm  |        +--- gauge 0-100 ---+     |
|  PIR:   ATTIVO           |        |        54         |     |
|  energia mov.:   72      |        |     "attivo"      |     |
|  energia still.: 30      |        +-------------------+     |
+--------------------------+----------------------------------+
|  C. GRAFICI                                                 |
|  C1: serie temporale energia moving/still + soglia           |
|  C2: barre per-gate (9 gate x moving/still) = "dov'e'"       |
|  C3: timeline presenza radar vs PIR (confronto visivo)       |
+-------------------------------------------------------------+
|  D. SESSIONE ED ESPORTAZIONE                                 |
|  scenario:[______] trial:[__] gt:[0|1]   <- metadati test    |
|  [Avvia] [Stop] [Scarica CSV] [Reset statistiche]            |
|  durata 12:34 | radar 97.2% | PIR 41.0% | eventi 3           |
+-------------------------------------------------------------+
\end{verbatim}
\end{minipage}
\caption{Struttura della pagina della dashboard. Il grafico C3 --- la timeline
affiancata di presenza radar e PIR --- è l'elemento che rende visivamente immediato
il risultato del \cref{cap:confronto}.}
\label{fig:webui-layout}
\end{figure*}

Due scelte di progetto meritano una nota.

\begin{itemize}[leftmargin=1.4em,itemsep=3pt]
  \item Il grafico C3 mostra \textbf{radar e PIR affiancati sulla stessa base dei
        tempi}. È l'elemento più efficace dal punto di vista dimostrativo: un
        osservatore vede la riga del radar continua e quella del PIR a impulsi
        isolati mentre una persona sta ferma davanti al sensore. È il risultato
        della tesi reso visibile in tempo reale.
  \item L'area D contiene i \textbf{metadati sperimentali} (scenario, trial, ground
        truth) con gli stessi nomi usati dalla catena seriale. In questo modo il CSV
        scaricato dal browser è direttamente utilizzabile dagli stessi script di
        analisi, senza conversioni.
\end{itemize}

\subsection{Statistiche di sessione}

Le statistiche sono calcolate \emph{nel browser} e non sull'ESP32, per non caricare
il microcontrollore con un compito che non gli è proprio:

\begin{itemize}[leftmargin=1.4em,itemsep=1pt]
  \item durata della sessione e percentuale di campioni con presenza, separatamente
        per radar e PIR;
  \item numero di eventi di attivazione (fronti $0 \to 1$) per ciascun sensore, che
        è la stima dei falsi positivi in tempo reale;
  \item distanza minima, media e massima; energia media;
  \item tempo trascorso dall'ultima rilevazione --- l'informazione più rilevante
        nello scenario di soccorso;
  \item vitalità: valore corrente, media mobile, minimo e massimo di sessione.
\end{itemize}

\section{Piano di realizzazione e criteri di accettazione}
\label{sec:webui-piano}

Lo sviluppo è organizzato in cinque passi, ciascuno con un criterio di accettazione
verificabile. La suddivisione serve a garantire che ogni passo produca qualcosa di
utilizzabile anche se i successivi non vengono completati.

\begin{enumerate}[leftmargin=1.8em,itemsep=3pt]
  \item \textbf{Server e pagina statica} --- l'ESP32 serve una pagina da LittleFS in
        entrambe le modalità WiFi. \emph{Accettazione}: la pagina si apre da browser
        collegandosi all'Access Point del nodo.
  \item \textbf{Flusso WebSocket} --- invio del messaggio di
        \cref{lst:ws} a 5~Hz. \emph{Accettazione}: i valori sulla pagina cambiano
        insieme al sensore, senza ritardo percepibile e senza accumulo di latenza su
        dieci minuti di funzionamento.
  \item \textbf{Visualizzazione} --- stato live, i tre grafici, la gauge di vitalità.
        \emph{Accettazione}: le barre per-gate indicano il gate corretto quando il
        soggetto si sposta di 0,75~m.
  \item \textbf{Registrazione ed esportazione} --- accumulo nel browser ed
        esportazione CSV. \emph{Accettazione}: il CSV scaricato dal browser e quello
        prodotto dalla catena seriale, per la stessa sessione, sono
        \textbf{equivalenti}; gli script di analisi producono le stesse metriche su
        entrambi. Questo è il test che chiude il requisito N2.
  \item \textbf{Indice di vitalità a bordo} --- porting dell'algoritmo del
        \cref{cap:vitalita}. \emph{Accettazione}: l'indice calcolato sull'ESP32 e
        quello calcolato in Python sullo stesso CSV coincidono entro l'errore di
        arrotondamento.
\end{enumerate}

\section{Riferimenti di progetto}

Come riferimento di interfaccia è stato esaminato il \emph{LD2410
Configurator}~\cite{ld2410configurator}, un configuratore web open source che dialoga
con il modulo tramite Web Serial e Web Bluetooth. Da esso sono stati ripresi due
elementi: la rappresentazione a barre delle energie per-gate con la soglia
sovrapposta, che è il modo più leggibile di mostrare venti canali insieme, e
l'impostazione ``una pagina, nessuna installazione''. Non sono state riprese le
funzioni di configurazione dei parametri, che nel nostro caso sarebbero controproducenti:
le soglie devono restare fisse per tutta la campagna (\cref{sec:soglie}), e un
pulsante che le modifica dalla dashboard è un rischio, non una funzione.

\paragraph{Avvertenza sulla fonte.} Il LD2410 Configurator è un progetto di comunità
e non un prodotto Hi-Link: vale come riferimento di interfaccia e di implementazione,
non come fonte di dati tecnici sul modulo. Per quelli si è fatto riferimento
esclusivamente ai documenti ufficiali~\cite{hilinkManualV104,hilinkProtoV107}.

\todo{completare con gli screenshot della dashboard realizzata e con l'esito
del test di accettazione del passo 4 (equivalenza dei due CSV)}
